\section{Canonical structure and the scope of approximation}
\label{sec:source-scope}

\subsection{The canonical score identity}

The word \emph{canonical} means that the coefficients come from the
unchanged complete theta density \eqref{eq:source}, with its equal
square-mode coefficients and full reflection identity. It does not
mean that every positive Jacobi matrix has the same Fourier geometry.
The full source enters through more information than the lower
score bound \eqref{eq:confinement}. At time zero put \(V=-\log\Phi\).
For the probability-normalized polynomials,
\begin{equation}
 a_n\int_\R V'(u)q_n(u)q_{n-1}(u)\,\dd\mu(u)=n,\qquad n\ge1.
 \label{eq:canonical-score}
\end{equation}
Indeed, integration of \((q_nq_{n-1}\Phi)'\) gives
\[
 \int V'q_nq_{n-1}\,\dd\mu
  =\int q_n'q_{n-1}\,\dd\mu+\int q_nq_{n-1}'\,\dd\mu=n/a_n.
\]
The second term vanishes by degree. If \(\kappa_n\) is the leading
coefficient of \(q_n\), the recurrence gives
\(\kappa_{n-1}=a_n\kappa_n\), so the first term is \(n/a_n\).
The tails justify the integration. In operator notation
\eqref{eq:canonical-score} reads
\(a_n[V'(J)]_{n,n-1}=n\), where the matrix element means the
displayed multiplication integral. It is not obtained by substituting
the finite matrix \(J_N\) into \(V'\); paths outside a finite block
are retained in the full multiplication operator.

The source is in fact characterized by all these identities under
the stated integrability hypothesis. This identifies where source
information lies, while keeping its distinction from a zero theorem.

\begin{proposition}[All-degree canonical-score characterization]
\label{prop:score-rigidity}
Let \(\nu\) be an even probability measure of infinite support on
\(\R\), and assume
\[
 \int_\R \e^{R|u|}(1+|V'(u)|)\,\dd\nu(u)<\infty
                      \qquad\hbox{for every }R>0.
\]
Construct its own orthonormal polynomials \(q_n^\nu\) and positive
Jacobi links \(a_n^\nu\). If
\begin{equation}
 a_n^\nu\int V'q_n^\nu q_{n-1}^\nu\,\dd\nu=n
                              \qquad(n\ge1),
 \label{eq:score-rigidity-hypothesis}
\end{equation}
then \(\nu=\mu\). Conversely, \(\mu\) satisfies these conditions.
\end{proposition}

\begin{proof}
Define \(\mathcal D(p)=\int(V'p-p')\,\dd\nu\).
For the measure's own polynomials, the same degree argument gives
\(\int(q_n^\nu q_{n-1}^\nu)'\,\dd\nu=n/a_n^\nu\).
Thus \(\mathcal D\) vanishes on each \(q_n^\nu q_{n-1}^\nu\).
These odd polynomials have successive odd degrees \(2n-1\) and
nonzero leading coefficients, so they form a triangular basis of
the odd polynomial space. On even polynomials \(\mathcal D\)
vanishes by parity, since \(V'\) is odd. Therefore
\(\mathcal D(p)=0\) for every polynomial.

Expand \(\e^{\ii tu}\) in its power series. The integrability
hypothesis dominates the series, its derivative, and its product
with \(V'\), giving
\[
 \int V'(u)\e^{\ii tu}\,\dd\nu(u)
     =\ii t\int\e^{\ii tu}\,\dd\nu(u)\qquad(t\in\R).
\]
Fourier inversion for a smooth compactly supported test function,
whose Fourier transform and its first polynomial multiple are
integrable, now gives
\(\int\varphi'\,\dd\nu=\int V'\varphi\,\dd\nu\).
In distributions this is \(\nu'=-V'\nu\). Multiplication by
the smooth function \(\e^V\) yields \((\e^V\nu)'=0\).
A locally finite measure with zero distributional derivative is a
constant multiple of Lebesgue measure: every compactly supported
smooth test of integral zero is the derivative of another such test.
Hence \(\dd\nu=c\e^{-V}\dd u=c\Phi\,\dd u\), and normalization
forces \(c=1/Z\).

For the converse, \eqref{eq:canonical-score} has already been proved.
The first positive summand bounds \(\Phi\) below by a constant times
\(\e^{|u|}\rho(u)^2\e^{-\rho(u)}\). Together with
\eqref{eq:tail} this gives \(|V'(u)|\le C\rho(u)\), so the complete
theta measure satisfies the required weighted exponential moments.
\end{proof}

The proposition uses all degrees. Its proof supplies neither a
quantitatively conditioned finite update nor a lower Fourier modulus.
The research question is whether this complete-source information
can produce a useful estimate of the type required in
Theorem~\ref{thm:contour} or Lemma~\ref{lem:zero-motion}.

\subsection{Scope of the general estimates}

The exact quadrature error law applies to any even probability
measure of infinite support with all absolute exponential moments.
The quantitative phase proof additionally uses the polynomial-weighted
cosh bound and the adjacent coefficient envelope. These premises
are not unique to theta.

For example, take a finite positive weighted-square sum
\(\sum_{m=1}^{M}c_ms_m(u)\) on \(u\ge0\) and its even extension.
Every \(c_m>0\) cancels from the corresponding summand score.
The same positive-mean argument proves \eqref{eq:score} and
\eqref{eq:confinement}, and all tail estimates needed for
Lemma~\ref{lem:general} hold. The even extension is locally absolutely
continuous even if higher odd derivatives do not match at zero.
In integration by parts for \(uP_n^2w\), the factor \(u\)
eliminates the junction contribution. Thus the same cosh moment,
coefficient envelope, complete error law and phase bound hold for
these sources. The known weighted-square collision examples
\cite{three,four} provide supporting witnesses to the limitation;
the numerical values of these controls do not enter the approximation proofs.

For orientation, the existing control comparisons have the following
distinct scopes. Euler-factor translates retain the generic exact
law, but the tested \(c=1/2,p=23,\lambda=0\) example fails the
adjacent envelope at index zero. Reflected super-Gaussian sources
have the exponential moments for that law but fail the exponential
score-confinement premise. Smooth shape and positive-primitive
controls lack the exact canonical summand formula; their weaker
confinement inequalities have not been tested here. Atoms allow the
law only before their monic norms vanish and lack a positive
full-line density; compact triangular extensions also lack that
global positivity. RH-conditional endpoint constructions retain
their hypothesis and do not supply unchanged theta arithmetic.
Changing the observation while keeping the source fails the literal
Fourier-readout premise, and relaxed profile measures do not identify
the full physical source. These are premise comparisons, not
universal exclusion theorems; the detailed catalog is cited in
the accompanying integration record.

A condition can hold near an isolated real zero of a control and
fail near a nonreal zero elsewhere. In particular, a successful
local disk comparison is not required to reject every control.
A global conclusion must account for all intended regions. If a
control has a nonreal zero inside a contour, the conjunction of
successful boundary separation and zero model count there must fail.
Failure of the canonical source formula alone is not proof that an
unspecified derived inequality fails.

\subsection{From source identification to zero geometry}

The steps established here are: source confinement bounds the
Jacobi coefficients; the coefficients determine Gaussian quadrature;
confluent interpolation expresses its complete error as
\(L_NE_N\); the source exponential moment makes \(E_N\) close to one;
and the directed enclosure yields pointwise exclusions or conditional
contour counts. The next missing estimate concerns the position of
the actual canonical readout relative to that enclosure.

For background, Van Assche \cite[\S1.1 and \S2.3]{vanassche}
explains orthogonal-polynomial recurrences and their weight
deformations; de Boor \cite[\S7 and \S9]{deboor} explains repeated-node
interpolation; Topping \cite[\S9.8--9.9]{topping} treats the argument
principle and Rouch\'e's theorem. Teschl
\cite[\S4.1]{teschl} supplies the spectral oscillation theory whose
variable must be distinguished from Fourier height in
\eqref{eq:two-zero-problems}. Romanov
\cite[\S3.2, Proposition~4]{romanov} exhibits a positive integral
identity leading to Hermite--Biehler geometry for a canonical
system. No faithful identification of our actual Fourier readouts
with such a system has been established here.

Two bounded questions remain useful. First, can the complete
source give lower slope or cluster-separation estimates with
enough degree dependence to apply a local zero-motion or count
argument? Second, can the canonical score identities yield a
relation for the readout and a justified companion function whose
phase obeys a usable comparison law? Either question needs an
actual quantitative conclusion and its own control comparison.
Neither is answered by the phase bound alone. The next sections
develop a local alternative: use a paid two-zero contour to form regular
cluster coordinates, preserve the signed physical error, and estimate
the deflated source curvature. The supporting certificate in
Appendix~\ref{sec:gt-source-certification} supplies one fixed tube.
A family covering growing heights, or a degree-propagated zero
confinement theorem, remains unestablished.
